\input{preamble}

\title{Section 3.3 --- Quadratic Functions}

\begin{document}
\maketitle
\section*{Solving Quadratic Equations}
For each quadratic equation:
\begin{itemize}
\item If it is in standard form, solve by factoring.
\item If it is in a form to allow it to be easily solved using the square root property, solve using the square root property.
\end{itemize}
\begin{smenum}
\mitemxxx
{$x^2-5x+6=0$}
{$x^2=9$}
{$(x-5)^2=49$}
\mitemxxx
{$4x^2+4x-15=0$}
{$3x^2-16x=12$}
{$2(x-3)^2+4=6$}
\end{smenum}
\section*{Graphing Quadratic Functions}
For each of the following quadratic functions:
\begin{clist}
\item Find the $x$- and $y$-intercepts for the graph.
\item Using your $x$-intercepts, find the equation of the axis of symmetry and the coordinates of the vertex.
\item Create a table of $x$ and $y$ values that include at least the intercepts, the point with the same $y$-coordinate as the $y$-intercept, and two integer points on each side of the vertex.
\item Sketch the axis of symmetry as a dashed line, and plot the vertex. (Estimate fractions if necessary.)
\item Graph the function in the space provided, plotting all appropriate points from your table.
\end{clist}
\begin{smenum}
\mitemxx
{$f(x)=x^2-6x+8$\\\GridSquareFitEO[10]}
{$g(x)=x^2-4x-5$\\\GridSquareFitEO[10]}
\clearpage
\mitemxx
{$h(x)=\frac{1}{2}x^2-3x$\\\GridSquareFitEO[10]}
{$k(x)=2x^2+10x+8$\\\GridSquareFitEO[10]}
\mitemxx
{$p(x)=-\frac{1}{3}x^2-x+6$\\\GridSquareFitEO[10]}
{$q(x)=-x^2+8x-7$\\\GridSquareFitEO[10]}
\end{smenum}
\end{document}
